\section{阿贝尔群的直和}\label{ch11:section-direct-sum}

本节我们只考虑阿贝尔群. 通常, 我们将这种群中的运算写作加法. 用 $0$ 表示群中的单位元, $-x$ 表示 $x$ 的逆元, $nx$ 表示 $n$ 个 $x$ 的和 $x+\cdots+x$.

设 $G$ 是一个阿贝尔群, $\{G_{\alpha}\}_{\alpha\in J}$ 为 $G$ 的子群的一个加标族. 若 $G$ 中的每一个元素 $x$ 可以表示为群族 $G_{\alpha}$ 中有限个成员之和, 则称群族 $G_{\alpha}$ 生成 (generate) $G$. 由于 $G$ 为阿贝尔群, 我们总可以经过适当分组, 使得每一组中的元素都在同一子群 $G_{\alpha}$ 中. 因此 $x$ 可以表示成以下形式:
\[
x=x_{\alpha_{1}}+\cdots+x_{\alpha_{n}},
\]
其中指标 $\{\alpha_{i}\}$ 是两两不同的. 从而, 我们经常将 $x$ 写成形式和 $x=\sum_{\alpha\in J}x_{\alpha}$, 其中当 $\alpha$ 不是 $\alpha_{1},\cdots,\alpha_{n}$ 中的某一个时, 约定 $x_{\alpha}=0$.

如果群族 $G_{\alpha}$ 生成 $G$, 则称 $G$ 为群族 $G_{\alpha}$ 的和 (sum), 记作 $G=\sum_{\alpha\in J}G_{\alpha}$, 而当指标集为有限指标集 $\{1,\cdots,n\}$ 时, 也记作 $G=G_{1}+\cdots+G_{n}$.

设群族 $G_{\alpha}$ 生成 $G$, 并且对于每一个 $x\in G$, $x$ 的表示 $x=\sum x_{\alpha}$ 是唯一的. 也就是说, 假定对于每一个 $x\in G$, 只有一个 $J$-串 $(x_{\alpha})_{\alpha\in J}$ 使得除有限个 $\alpha$ 外都有 $x_{\alpha}=0$, 并且 $x=\sum x_{\alpha}$. 这时 $G$ 称为群族 $G_{\alpha}$ 的直和 (direct sum), 记作
\[
G=\bigoplus_{\alpha\in J}G_{\alpha},
\]
当指标集有限时, 也记作 $G=G_{1}\oplus\cdots\oplus G_{n}$.

\begin{example}\label{ch11:ex-direct-sum-r-infty}
笛卡儿积 $\mathbb{R}^{\omega}$ 关于按坐标相加的运算构成阿贝尔群. 当 $i\neq n$ 时, 使得 $x_{i}=0$ 的所有元素串 $(x_{i})$ 所构成的集合 $G_{n}$ 是一个同构于 $\mathbb{R}$ 的子群. 群族 $G_{n}$ 生成 $\mathbb{R}^{\omega}$ 的子群 $\mathbb{R}^{\infty}$, 事实上, $\mathbb{R}^{\infty}$ 是这些群的直和.
\end{example}

下面的引理指出了直和的一个有用的特征, 称为直和的扩展条件 (extension condition).

\begin{lemma}\label{ch11:lemma-extension-condition}
设 $G$ 是一个阿贝尔群, $\{G_{\alpha}\}$ 为 $G$ 的子群的一个族. 若 $G$ 为群族 $G_{\alpha}$ 的直和, 则 $G$ 满足以下条件:

$(*)$ 对于任何阿贝尔群 $H$ 以及任何同态 $h_{\alpha}\colon G_{\alpha}\to H$ 的族, 存在一个同态 $h\colon G\to H$, 使得对于每一个 $\alpha$, $h$ 在 $G_{\alpha}$ 上的限制等于 $h_{\alpha}$.

此外, $h$ 还是唯一的. 反之, 如果群族 $G_{\alpha}$ 生成 $G$ 并且扩展条件 $(*)$ 成立, 则 $G$ 为群族 $G_{\alpha}$ 的直和.
\end{lemma}

\begin{proof}
我们首先证明: 如果 $G$ 具有上述扩展性质, 则 $G$ 为群族 $G_{\alpha}$ 的直和. 假设 $x=\sum x_{\alpha}=\sum y_{\alpha}$, 我们来证明对任何指标 $\beta$ 有 $x_{\beta}=y_{\beta}$. 令 $H$ 表示群 $G_{\beta}$, 设 $h_{\alpha}\colon G_{\alpha}\to H$ 满足以下条件: 当 $\alpha\neq\beta$ 时为平凡同态, 当 $\alpha=\beta$ 时为恒等同态. 设 $h\colon G\to H$ 为由假设给出的同态族 $h_{\alpha}$ 的扩充, 则
\[
h(x)=\sum h_{\alpha}(x_{\alpha})=x_{\beta},\qquad
h(x)=\sum h_{\alpha}(y_{\alpha})=y_{\beta}.
\]
于是 $x_{\beta}=y_{\beta}$.

以下证明: 如果 $G$ 为群族 $G_{\alpha}$ 的直和, 则扩展条件成立. 对于给定的同态 $h_{\alpha}$, 我们定义 $h(x)$ 如下: 若 $x=\sum x_{\alpha}$, 令 $h(x)=\sum h_{\alpha}(x_{\alpha})$. 由于此处为有限和, 所以上述定义有意义. 由于 $x$ 的表示是唯一的, 所以 $h$ 的定义是确切的. 易于验证, $h$ 即为所求的同态. 由于对于每一个 $\alpha$, $h$ 在 $G_{\alpha}$ 上的限制等于 $h_{\alpha}$, $h$ 必定满足上述等式, 因此 $h$ 是唯一的.
\end{proof}

这个引理使得有关直和的许多结果易于证明.

\begin{corollary}\label{ch11:cor-associativity}
设 $G=G_{1}\oplus G_{2}$. 假定 $G_{1}$ 为子群族 $\{H_{\alpha}\}_{\alpha\in J}$ 的直和, $G_{2}$ 为子群族 $\{H_{\beta}\}_{\beta\in K}$ 的直和, 其中指标集 $J$ 与 $K$ 无交. 则 $G$ 为子群族 $\{H_{\gamma}\}_{\gamma\in J\cup K}$ 的直和.
\end{corollary}

\begin{proof}
若 $h_{\alpha}\colon H_{\alpha}\to H$ 和 $h_{\beta}\colon H_{\beta}\to H$ 皆为同态族, 根据上述引理, 它们的扩充同态为 $h_{1}\colon G_{1}\to H$ 和 $h_{2}\colon G_{2}\to H$. 从而 $h_{1}$ 和 $h_{2}$ 可以扩充为同态 $h\colon G\to H$.
\end{proof}

作为例子, 上述推论蕴涵着下列事实:
\[
(G_{1}\oplus G_{2})\oplus G_{3}=G_{1}\oplus G_{2}\oplus G_{3}=G_{1}\oplus(G_{2}\oplus G_{3}).
\]

\begin{corollary}\label{ch11:cor-quotient}
若 $G=G_{1}\oplus G_{2}$, 则 $G/G_{2}$ 同构于 $G_{1}$.
\end{corollary}

\begin{proof}
设 $H=G_{1}$, $h_{1}\colon G_{1}\to H$ 为恒等同态, $h_{2}\colon G_{2}\to H$ 为平凡同态. 假定 $h\colon G\to H$ 为 $G$ 的扩充. 则 $h$ 是以 $G_{2}$ 为核的满射.
\end{proof}

在许多场合, 人们希望对给定的阿贝尔群的一个族 $\{G_{\alpha}\}$ 找到一个群 $G$ 使得它含有同构于 $G_{\alpha}$ 的子群 $G'_{\alpha}$, 并且 $G$ 还是这些子群的直和. 事实上, 这总是可以办到的. 为此, 我们需要引入外直和这个概念.

\begin{definition}[external direct sum]\label{ch11:def-external-direct-sum}
设 $\{G_{\alpha}\}_{\alpha\in J}$ 为阿贝尔群的一个加标族. 假定 $G$ 是一个阿贝尔群, $i_{\alpha}\colon G_{\alpha}\to G$ 是单同态的一个族, 使得 $G$ 为群族 $i_{\alpha}(G_{\alpha})$ 的直和, 那么称 $G$ 为群族 $G_{\alpha}$ 的关于单同态族 $i_{\alpha}$ 的外直和 (external direct sum).
\end{definition}

当然, 群 $G$ 并不唯一. 然而我们将证明在不区别同构的群的意义下它是唯一的. 这里有一个构造 $G$ 的办法.

\begin{theorem}\label{ch11:thm-external-existence}
对于给定的阿贝尔群的一个族 $\{G_{\alpha}\}_{\alpha\in J}$, 存在一个阿贝尔群 $G$ 和单同态的一个族 $i_{\alpha}\colon G_{\alpha}\to G$, 使得 $G$ 为群族 $i_{\alpha}(G_{\alpha})$ 的直和.
\end{theorem}

\begin{proof}
首先考虑笛卡儿积
\[
\prod_{\alpha\in J}G_{\alpha}.
\]
如果我们通过按坐标相加来定义两个 $J$-元素串的加法, 那么它是一个阿贝尔群. 设 $G$ 是一个以 $(x_{\alpha})_{\alpha\in J}$ 作为其元素的笛卡儿积的子群, 其中 $(x_{\alpha})_{\alpha\in J}$ 满足以下条件: 除去有限多个 $\alpha$ 之外, $x_{\alpha}=0_{\alpha}$, $0_{\alpha}$ 是 $G_{\alpha}$ 的单位元. 对于给定的指标 $\beta$, 定义 $i_{\beta}\colon G_{\beta}\to G$ 使得 $i_{\beta}(x)$ 以 $x$ 为其第 $\beta$ 个坐标, 并且对于任何 $\alpha\neq\beta$ 以 $0_{\alpha}$ 为其第 $\alpha$ 个坐标. 易见, $i_{\beta}$ 为单同态. 由于 $G$ 的每一个元素 $x$ 仅有有限个坐标不是单位元, 也易见, $G$ 中的每一个元素 $x$ 可以唯一地表示成群族 $i_{\beta}(G_{\beta})$ 中的元素的有限和.
\end{proof}

可以直接将刻画直和的扩展条件转化为刻画外直和的扩展条件.

\begin{theorem}\label{ch11:thm-external-extension}
设 $\{G_{\alpha}\}_{\alpha\in J}$ 为阿贝尔群的一个加标族, $G$ 是一个阿贝尔群, $i_{\alpha}\colon G_{\alpha}\to G$ 为同态的一个族. 如果每一个 $i_{\alpha}$ 都是单同态, 并且 $G$ 是群族 $i_{\alpha}(G_{\alpha})$ 的直和, 则 $G$ 满足以下扩展条件:

$(*)$ 对于任何阿贝尔群 $H$ 以及任何同态族 $h_{\alpha}\colon G_{\alpha}\to H$, 存在一个同态 $h\colon G\to H$, 使得对于每一个 $\alpha$ 都有 $h\circ i_{\alpha}=h_{\alpha}$.

此外, $h$ 还是唯一的. 反之, 如果群族 $i_{\alpha}(G_{\alpha})$ 生成 $G$ 并且扩展条件 $(*)$ 成立, 则每一个 $i_{\alpha}$ 都是单同态, 并且 $G$ 为群族 $i_{\alpha}(G_{\alpha})$ 的直和.
\end{theorem}

\begin{proof}
只要证明: 如果扩展条件成立, 则每一个 $i_{\alpha}$ 为单同态. 对于任意给定的指标 $\beta$, 令 $H=G_{\beta}$ 以及 $h_{\alpha}\colon G_{\alpha}\to H$ 满足以下条件: 当 $\alpha=\beta$ 时为恒等同态, 当 $\alpha\neq\beta$ 时为平凡同态. 令 $h\colon G\to H$ 表示满足假设条件的扩充, 则特别有 $h\circ i_{\beta}=h_{\beta}$, 这蕴涵着 $i_{\beta}$ 为单射.
\end{proof}

一个直接的推论便是直和的唯一性定理.

\begin{theorem}[uniqueness of direct sum]\label{ch11:thm-direct-sum-uniqueness}
设 $\{G_{\alpha}\}_{\alpha\in J}$ 为阿贝尔群的一个族. 设 $G$ 和 $G'$ 都是阿贝尔群, $i_{\alpha}\colon G_{\alpha}\to G$ 和 $i'_{\alpha}\colon G_{\alpha}\to G'$ 都是单同态族, 其中 $G$ 是群族 $i_{\alpha}(G_{\alpha})$ 的直和, $G'$ 是群族 $i'_{\alpha}(G_{\alpha})$ 的直和. 则存在一个唯一同构 $\phi\colon G\to G'$ 使得 $\phi\circ i_{\alpha}=i'_{\alpha}$ 对于每一个 $\alpha$ 成立.
\end{theorem}

\begin{proof}
应用前面的引理. 由于 $G$ 为群族 $G_{\alpha}$ 的外直和, $\{i'_{\alpha}\}$ 是同态的一个族, 所以存在一个唯一同态 $\phi\colon G\to G'$ 使得 $\phi\circ i_{\alpha}=i'_{\alpha}$ 对于每一个 $\alpha$ 成立. 类似地, 由于 $G'$ 是群族 $G_{\alpha}$ 的外直和, 并且 $\{i_{\alpha}\}$ 是同态的一个族, 从而存在一个唯一同态 $\psi\colon G'\to G$ 使得 $\psi\circ i'_{\alpha}=i_{\alpha}$ 对于每一个 $\alpha$ 成立. 于是 $\psi\circ\phi\colon G\to G$ 使得 $\psi\circ\phi\circ i_{\alpha}=i_{\alpha}$ 对于每一个 $\alpha$ 成立. 由于 $G$ 的恒等映射具有如上性质, 根据引理的唯一性部分得知: $\psi\circ\phi$ 必为 $G$ 上的恒等映射. 同理, $\phi\circ\psi$ 也必为 $G'$ 上的恒等映射.
\end{proof}

如果 $G$ 为相对于单同态族 $i_{\alpha}$ 的群族 $G_{\alpha}$ 的外直和, 我们有时也将 $G$ 写成 $G=\bigoplus G_{\alpha}$, 尽管群族 $G_{\alpha}$ 不是 $G$ 的子群. 也就是说, 我们等同群 $G_{\alpha}$ 与其在 $i_{\alpha}$ 下的像, 然后将 $G$ 视为通常的直和, 而不是外直和. 具体情况可依照上下文而定.

以下我们来讨论自由阿贝尔群.

\begin{definition}[free Abelian group]\label{ch11:def-free-abelian}
设 $G$ 是一个阿贝尔群, $\{a_{\alpha}\}$ 是 $G$ 中元素的一个加标族, $G_{\alpha}$ 是由 $a_{\alpha}$ 生成的 $G$ 的子群. 如果群族 $G_{\alpha}$ 生成 $G$, 我们也称元素族 $a_{\alpha}$ 生成 $G$. 如果每一个群 $G_{\alpha}$ 都是无限循环群, 并且 $G$ 为群族 $G_{\alpha}$ 的直和, 则称 $G$ 为以元素族 $\{a_{\alpha}\}$ 为基 (base) 的自由阿贝尔群 (free Abelian group).
\end{definition}

直和的扩展条件蕴涵着以下自由阿贝尔群的扩展条件.

\begin{lemma}\label{ch11:lemma-free-abelian-extension}
设 $G$ 是一个阿贝尔群, $\{a_{\alpha}\}_{\alpha\in J}$ 为 $G$ 中元素的一个族, 它们生成 $G$. 则 $G$ 是一个以 $\{a_{\alpha}\}$ 为基的自由阿贝尔群的充分必要条件是, 对于任何一个阿贝尔群 $H$ 以及 $H$ 中任意元素的一个族 $\{y_{\alpha}\}$, 存在一个 $G$ 到 $H$ 的同态 $h$ 使得 $h(a_{\alpha})=y_{\alpha}$ 对于每一个 $\alpha$ 成立. 这时, $h$ 是唯一的.
\end{lemma}

\begin{proof}
令 $G_{\alpha}$ 表示由 $a_{\alpha}$ 生成的 $G$ 的子群. 首先假定扩展性质成立. 我们来证明每一个 $G_{\alpha}$ 都是无限循环群. 假设有某个指标 $\beta$ 使得元素 $a_{\beta}$ 生成 $G$ 的一个有限循环子群. 若取 $H=\mathbb{Z}$, 则任何同态 $h\colon G\to H$ 都不会满足将每一个 $a_{\alpha}$ 映射到 $1$. 因为 $a_{\beta}$ 是有限阶的, 而 $1$ 不是有限阶的! 要证明 $G$ 为群族 $G_{\alpha}$ 的直和, 只需应用引理 67.1 便可.

反之, 如果 $G$ 是以 $\{a_{\alpha}\}$ 为基的自由阿贝尔群, 则对于给定的 $H$ 中的元素族 $\{y_{\alpha}\}$, 存在同态 $h_{\alpha}\colon G_{\alpha}\to H$ 使得 $h_{\alpha}(a_{\alpha})=y_{\alpha}$ (因为 $G_{\alpha}$ 是无限循环群). 然后应用引理 67.1.
\end{proof}

\begin{theorem}\label{ch11:thm-rank-well-defined}
设 $G$ 是以 $\{a_{1},\cdots,a_{n}\}$ 为基的一个自由阿贝尔群, 则 $n$ 是由 $G$ 唯一确定的.
\end{theorem}

\begin{proof}
群 $G$ 同构于 $n$-重积 $\mathbb{Z}\times\cdots\times\mathbb{Z}$, 子群 $2G$ 对应于积 $(2\mathbb{Z})\times\cdots\times(2\mathbb{Z})$. 这时商群 $G/2G$ 与集合 $(\mathbb{Z}/2\mathbb{Z})\times\cdots\times(\mathbb{Z}/2\mathbb{Z})$ 一一对应, 从而 $G/2G$ 的基数为 $2^{n}$, 于是 $n$ 由 $G$ 唯一确定.
\end{proof}

如果 $G$ 为具有有限基的自由阿贝尔群, 则 $G$ 的基元素的个数叫做 $G$ 的秩 (rank).

\begin{exercise}\label{ch11:ex-direct-sum-1}
设 $G=\sum G_{\alpha}$. 证明: 这个和式为直和当且仅当等式
\[
x_{\alpha_{1}}+\cdots+x_{\alpha_{n}}=0
\]
蕴涵着每一个 $x_{\alpha_{i}}=0$. (这里 $x_{\alpha_{i}}\in G_{\alpha_{i}}$ 且诸指标 $\alpha_{i}$ 两两不同.)
\end{exercise}

\begin{exercise}\label{ch11:ex-direct-sum-2}
证明: 若 $G_{1}$ 为 $G$ 的一个子群, 那么满足 $G=G_{1}\oplus G_{2}$ 的 $G$ 的子群 $G_{2}$ 可能不存在. [提示: 取 $G=\mathbb{Z}$ 及 $G_{1}=2\mathbb{Z}$.]
\end{exercise}

\begin{exercise}\label{ch11:ex-direct-sum-3}
设 $G$ 是以 $\{x,y\}$ 为基的自由阿贝尔群, 证明 $\{2x+3y,x-y\}$ 也是 $G$ 的基.
\end{exercise}

\begin{exercise}\label{ch11:ex-direct-sum-4}
阿贝尔群中元素 $a$ 的阶 (order) 是指满足条件 $ma=0$ 的最小正整数 $m$, 若这样的正整数不存在, 则称 $a$ 的阶为无穷大. $a$ 的阶便是它所生成的子群的阶.

\begin{enumerate}[(a)]
\item 证明: $G$ 的所有有限阶元素构成 $G$ 的一个子群, 称为挠子群 (torsion subgroup).
\item 证明: 若 $G$ 是自由阿贝尔群, 则它没有有限阶元素.
\item 证明: 有理数加群没有有限阶元素, 但它不是自由阿贝尔群. [提示: 若 $\{a_{\alpha}\}$ 是一个基, 用这个基来表示 $\frac{1}{2}a_{\alpha}$.]
\end{enumerate}
\end{exercise}

\begin{exercise}\label{ch11:ex-direct-sum-5}
举例说明: 一个秩为 $n$ 的自由阿贝尔群 $G$ 具有秩为 $n$ 的子群 $H$, $H\neq G$.
\end{exercise}

\begin{exercise}\label{ch11:ex-direct-sum-6}
证明以下结论:

\textbf{定理.} 若 $A$ 是秩为 $n$ 的自由阿贝尔群, 则 $A$ 的任何子群 $B$ 都是秩不超过 $n$ 的自由阿贝尔群.

证明: 可以假设 $A=\mathbb{Z}^{n}$, 即 $\mathbb{Z}$ 的 $n$-重笛卡儿积. 设 $\pi_{i}\colon\mathbb{Z}^{n}\to\mathbb{Z}$ 是到第 $i$ 个坐标的投影. 任意给定 $m\leqslant n$, 设 $B_{m}$ 为 $B$ 的子集, 其元素 $x$ 满足条件: 当 $i>m$ 时, $\pi_{i}(x)=0$, 则 $B_{m}$ 为 $B$ 的子群. 考虑 $\mathbb{Z}$ 的子群 $\pi_{m}(B_{m})$. 如果它是非平凡的, 选取 $x_{m}\in B_{m}$ 使得 $\pi_{m}(x_{m})$ 为这个子群的生成元. 否则, 设 $x_{m}=0$.

\begin{enumerate}[(a)]
\item 证明: 对于每一个 $m$, $\{x_{1},\cdots,x_{m}\}$ 生成 $B_{m}$.
\item 证明: 对于每一个 $m$, $\{x_{1},\cdots,x_{m}\}$ 所有非零元素构成 $B_{m}$ 的一个基.
\item 证明: $B_{n}=B$ 是秩不超过 $n$ 的自由阿贝尔群.
\end{enumerate}
\end{exercise}